\section{Glory in Captivity}

At Shiloh the word of the Lord is rare, Eli's sight is failing, and yet the lamp
of God has not gone out.  Near the Ark, the child Samuel hears his name and
learns to answer (1 Sam 3:1--10).  The sanctuary is not abandoned by grace, but
its priestly house stands under judgment.  The sons of Eli abuse their office;
the people, facing Philistine power, will soon repeat the same error in another
form.  They will treat the Ark not as the sign of a covenant to obey, but as a
force to deploy.

\subsection{The shout before the defeat}

Israel camps at Ebenezer; the Philistines camp at Aphek.  After Israel's first
defeat, the elders send to Shiloh for the Ark, reasoning that its arrival will
save them from their enemies.  Hophni and Phinehas come with it.  When the Ark
enters the camp, Israel shouts until the earth rings.  The Philistines hear,
remember the mighty deeds of Israel's God in Egypt, and brace themselves in
fear (1 Sam 4:1--9).

Then Israel is defeated.  Hophni and Phinehas die.  The Ark is taken
(1 Sam 4:10--11).  Aphek remains the Philistine camp, not an Ark station; the
captured Ark leaves the Israelite camp at Ebenezer (1 Sam 5:1).

A messenger runs to Shiloh with torn clothes and earth upon his head.  Eli can
bear the report of military disaster and the death of his sons, but when he
hears that the Ark has been captured, he falls backward and dies.  His
daughter-in-law gives birth amid the catastrophe and names the child Ichabod:
the glory has departed from Israel (1 Sam 4:12--22).  Her lament names the
felt disaster.  The next chapter reveals the theological reversal.  The Ark
has been captured; the Lord has not.

\subsection{Dagon on the threshold}

The Philistines carry the Ark from Ebenezer to Ashdod and set it beside Dagon
in his temple.  In the morning Dagon lies face down before the Ark.  They put
him back.  On the next morning he lies fallen again, his head and hands cut off
at the threshold (1 Sam 5:1--5).  No Israelite army enters the scene.  The
supposed trophy has become the silent center of an iconoclastic victory.

The hand of the Lord lies heavy on Ashdod.  The Ark is sent to Gath; affliction
follows.  It is sent onward, and the final city cries out before its arrival
(1 Sam 5:6--12).  Here the ancient witnesses preserve a real route difference.
The Hebrew, Vulgate, and Douay line names Ekron as the third stop; the
registered Brenton Septuagint names Ashkelon.  The Hebrew/Vulgate itinerary may
therefore be told as Ashdod--Gath--Ekron, while the Greek branch must remain
visible rather than being silently corrected.  Gaza appears later among the
five Philistine lordships represented by guilt offerings; it is never a
narrated Ark stop (1 Sam 6:17).

For seven months the Ark remains in Philistine territory---seven months total,
not seven months in each city (1 Sam 6:1).  Priests and diviners advise a return
with reparation.  A new cart is prepared.  Two milk cows that have never borne
a yoke are harnessed; their calves are shut away.  If the cows take the road to
Beth-shemesh, the Philistines will recognize the Lord's hand (1 Sam 6:2--9).

\subsection{The straight road home}

The cows go straight toward Beth-shemesh, lowing as they walk and turning
neither right nor left.  The Philistine rulers follow as far as the border.
In the valley the people are reaping wheat.  They lift their eyes, see the Ark,
and rejoice.  The cart stops in the field of Joshua beside a great stone; its
wood becomes fuel, and the cows are offered to the Lord
(1 Sam 6:10--18).

Joy does not abolish holiness.  The narrative records a deadly judgment at
Beth-shemesh, and the people mourn and ask who can stand before the holy Lord
(1 Sam 6:19--20).  The verse's wording and casualty figure are exceptionally
unstable across textual and translation traditions: witnesses variously speak
of seeing, looking upon, or looking into the Ark, and they do not agree on the
number struck.  A thrilling retelling has no need to pretend certainty where
the text does not give it.

Messengers go to Kiriath-jearim.  Its men come down, bring the Ark to the house
of Abinadab on the hill, and consecrate Eleazar to guard it
(1 Sam 6:21--7:1).  Twenty years pass before the narrative turns to Israel's
lament and return under Samuel (1 Sam 7:2--3).  That notice does not measure
the Ark's entire residence there; it remains until David comes to take it up
toward Zion (2 Sam 6:2; 1 Chr 13:5--6).

The whole circuit is a judgment on religious possession.  Israel cannot turn
the Ark into a weapon.  Philistia cannot turn it into a trophy.  Dagon cannot
stand before it.  Even the joyful people of Beth-shemesh cannot erase the
boundary of holiness.  Yet the Lord returns the Ark without an Israelite
rescue, guiding the untrained team along a straight road home.

\begin{studybox}{The captive Ark and the unconquered Lord}
Read 1 Samuel 4:21--22 beside 5:1--5.  ``Glory has departed'' voices Israel's
bereavement; Dagon's fall reveals that God's sovereignty has not departed.
Faith may grieve a real loss without imagining that God has been defeated.
The sign can be taken.  The Lord remains free.
\end{studybox}
